%==============================================================================
\appsection{Prior Matching with 3 Moments: The Skew-Normal Distribution}
\label{app:skew_normal}
%==============================================================================

This appendix demonstrates the law of total cumulance for a hierarchical model
where the unconditional distribution is Skew-Normal, requiring matching of
three cumulants rather than two.

\subsection{Law of Total Cumulance for Skewness}

The law of total cumulance provides a general decomposition for the $n$-th
cumulant. For the third cumulant $\kappa_3$ (equal to the third central moment,
related to skewness):
\begin{equation}
  \label{eq:total_skewness}
  \kappa_3(Y) = \E{\kappa_3(Y|X)} + 3\Cov{\E{Y|X}}{\Var{Y|X}} + \kappa_3(\E{Y|X})
\end{equation}
\noindent The three terms have intuitive interpretations:
\begin{enumerate}
  \item $\E{\kappa_3(Y|X)}$: Expected within-group skewness
  \item $3\Cov{\E{Y|X}}{\Var{Y|X}}$: Covariance between group means and variances
  \item $\kappa_3(\E{Y|X})$: Skewness of the group means
\end{enumerate}

\subsection{The Skew-Normal Distribution}

The Skew-Normal distribution has density:
\begin{equation}
  f(x; \xi, \omega, \alpha) = \frac{2}{\omega} \phi\left(\frac{x-\xi}{\omega}\right) \Phi\left(\alpha\frac{x-\xi}{\omega}\right)
\end{equation}
where $\xi$ is location, $\omega > 0$ is scale, $\alpha$ is the shape parameter
controlling skewness, and $\phi$, $\Phi$ are the standard normal PDF and CDF.

The cumulants are:
\begin{align}
  \kappa_1 &= \xi + \omega \delta \sqrt{\frac{2}{\pi}}, \quad \text{where } \delta = \frac{\alpha}{\sqrt{1+\alpha^2}} \\
  \kappa_2 &= \omega^2 \left(1 - \frac{2\delta^2}{\pi}\right) \\
  \kappa_3 &= \frac{4-\pi}{2} \left(\delta\sqrt{\frac{2}{\pi}}\right)^3 \omega^3
\end{align}

\subsection{Hierarchical Decomposition}

Suppose we wish to specify hyperpriors such that the unconditional distribution
of $\theta_i$ matches a direct prior $\text{Skew-Normal}(\xi^*, \omega^*, \alpha^*)$.
With three parameters, we must match three cumulants. The matching conditions
couple the hyperpriors for $\xi$, $\omega$, and $\alpha$ in a more complex
manner than the Normal case.

\subsection{Hyperprior Specification}

Consider the hyperprior structure:
\begin{align}
  \xi &\sim \text{Normal}(\mu_\xi, \tau_\xi^2) \\
  \omega^2 &\sim \text{Inverse-Gamma}(\alpha_\omega, \beta_\omega) \\
  \alpha &\sim \text{Normal}(\mu_\alpha, \tau_\alpha^2)
\end{align}
This introduces 6 hyperprior parameters. The 3 cumulant matching conditions
determine 3, leaving 3 degrees of freedom:
\begin{enumerate}
  \item $\rho \in (0,1)$: Variance partition ratio (same interpretation as Normal case).
  \item $\alpha_\omega > 2$: Scale concentration (analogous to $\alpha_0$).
  \item $\tau_\alpha^2 > 0$: Shape uncertainty (how much the skewness can vary
    across the population).
\end{enumerate}

\subsection{Numerical Solution}

The nonlinear dependence on $\delta(\alpha) = \alpha/\sqrt{1+\alpha^2}$ generally
precludes closed-form solutions. Given researcher choices, the system
reduces to 3 equations in 3 unknowns $(\mu_\xi, \beta_\omega, \mu_\alpha)$,
solvable via root-finding with Monte Carlo integration over the hyperprior
distributions to compute the implied cumulants.


\pagebreak
